\documentclass{ximera}

\title{Exercises for Functions and their Representations} 
\license{CC BY-NC-SA 4.0}

\begin{document}

\begin{abstract}
\end{abstract}
\maketitle

\begin{question}

In Exercises \ref{mappingfirst} - \ref{mappinglast}, determine whether or not the mapping diagram represents a function. Explain your reasoning. If the mapping does represent a function, state the domain, range, and represent the function as a set of ordered pairs.

\begin{problem}\label{mappingfirst}

\begin{tikzpicture}[>=Stealth]

% Left labels
\node at (-4.25,5) {Tennant};
\node at (-4.25,3) {Smith};
\node at (-4.25,1) {Calpadi};

% Center label
\node at (0,6) {$M$};

% Right labels
\node at (3.5,5) {Eleven};
\node at (3.5,3) {Twelve};
\node at (3.5,1) {Thirteen};
\node at (3.5,-1) {Fourteen};

% Arrows
\draw[->] (-2.5,5) -- (2.5,5);
\draw[->] (-2.5,5) -- (2.5,3);
\draw[->] (-2.5,3) -- (2.5,1);
\draw[->] (-2.5,1) -- (2.5,-1);

\end{tikzpicture}
\end{problem} 



\begin{problem}\label{mappinglast}

\begin{tikzpicture}[>=Stealth]

% Left labels
\node at (-4.25,5) {Hartnell};
\node at (-4.25,3) {Cushing};
\node at (-4.25,1) {Hurndall};
\node at (-4.25,-1) {Troughton};

% Center label
\node at (0,6) {$C$};

% Right labels
\node at (3.5,5) {One};
\node at (3.5,3) {Two};
\node at (3.5,1) {Three};

% Arrows
\draw[->] (-2.5,5) -- (2.5,5);
\draw[->] (-2.5,3) -- (2.5,5);
\draw[->] (-2.5,1) -- (2.5,5);
\draw[->] (-2.5,-1) -- (2.5,3);

\end{tikzpicture}
\end{problem} 
    
\end{question}

\begin{question}
    
In Exercises \ref{tablefirst} - \ref{tablelast}, determine whether or not the data in the given table represents $y$ as a function of $x$.  Explain your reasoning.  If the mapping does represent a function, state the domain, range, and represent the function as a set of ordered pairs.


\begin{problem}

\label{tablefirst}
    \[\begin{array}{|r||r|}  \hline

x  &  y  \\ \hline
 -3 &  3 \\  \hline
 -2 & 2  \\  \hline
  -1 &  1  \\  \hline
 0 &  0 \\  \hline
 1 & 1  \\  \hline
 2 &  2 \\  \hline
 3 & 3  \\  \hline

\end{array}\]
\end{problem}

\begin{problem}\label{tablelast}

    \[\begin{array}{|r||r|}  \hline

 x  & y  \\ \hline

 0 & 0 \\  \hline
 1 & 1  \\  \hline
 1 & -1  \\  \hline
 2 &  2 \\  \hline
 2 & -2  \\  \hline
 3 &  3 \\  \hline
 3 & -3  \\  \hline

\end{array}\]
\end{problem}

\end{question}




\begin{problem} 
Suppose $W$ is the set of words in the English language and we set up a mapping from $W$ into the set of natural numbers $\mathbb{N}$ as follows: $\text{word} \rightarrow \text{number of letters in the word}$.  Explain why this mapping is a function.  What would you need to know to determine the range of the function?
\end{problem}

\begin{problem}
    Suppose $L$ is the set of last names of all the people who have served or are currently serving as the President of the United States.   Consider the mapping from $L$ into $\mathbb{N}$ as follows:  $\text{last name} \rightarrow \text{number of their presidency}$.  For example,  $\text{Washington} \rightarrow 1$ and $\text{Obama} \rightarrow 44$.  Is this mapping a function?  What if we use full names instead of just last names?
    \begin{hint}
    Research what Grover Cleveland and Donald Trump have in common.
    \end{hint}
\end{problem} 

\begin{problem}
  Under what conditions would the time of day be a function of the outdoor temperature?   
\end{problem}

\begin{question}

For the functions $f$ described in Exercises \ref{buildfunctionfirst} - \ref{buildfunctionlast}, find $f(2)$, and find and simplify an expression for $f(x)$ that takes a real number $x$ and performs the following three steps in the order given.


\begin{problem}\label{buildfunctionfirst}
(1) multiply by 2; (2) add 3; (3) divide by 4. 

    
\end{problem}


\begin{problem}
(1) add 3; (2) multiply by 2; (3) divide by 4. 

\end{problem}

\begin{problem}
(1) divide by 4; (2) add 3; (3) multiply by 2.

\end{problem}


\begin{problem}
(1) multiply by 2; (2) add 3; (3) take the square root.

    Find $f(2)$, and find and simplify an expression for $f(x)$.
\end{problem}


\begin{problem}
(1) add 3; (2) multiply by 2; (3) take the square root.

    Find $f(2)$, and find and simplify an expression for $f(x)$.
\end{problem}

\begin{problem}\label{buildfunctionlast}
(1) add 3; (2) take the square root; (3) multiply by 2.

    Find $f(2)$, and find and simplify an expression for $f(x)$.
\end{problem}
    
\end{question}


\begin{question}
In Exercises \ref{funcnotationbasicfirst} - \ref{funcnotationbasiclast}, use the given function $f$ to find and simplify the following:

\begin{itemize}
\item $f(3)$
\item $f(-1)$
\item $f\left(\frac{3}{2} \right)$
\item  $f(4x)$
\item $4f(x)$
\item $f(-x)$
\item  $f(x-4)$
\item $f(x) - 4$
\item  $f\left(x^2\right)$
\end{itemize}


\begin{problem}\label{funcnotationbasicfirst} 
$f(x) = 2x+1$

\begin{itemize}
\item $f(3) = \answer{7}$
\item $f(-1) = \answer{-1}$
\item $f\left(\frac{3}{2} \right) = \answer{4}$
\item  $f(4x) = \answer{8x+1}$
\item $4f(x) = \answer{8x+4}$
\item $f(-x) = \answer{-2x+1}$
\item  $f(x-4) = \answer{2x-7}$
\item $f(x) - 4 = \answer{2x-3}$
\item  $f\left(x^2\right) = \answer{2x^2+1}$
\end{itemize}

\end{problem}  

\begin{problem}
$f(x) = 3 - 4x$

\begin{itemize}
\item $f(3) = \answer{-9}$
\item $f(-1) = \answer{7}$
\item $f\left(\frac{3}{2} \right) = \answer{-3}$
\item  $f(4x) = \answer{3-16x}$
\item $4f(x) = \answer{12-16x}$
\item $f(-x) = \answer{4x+3}$
\item  $f(x-4) = \answer{19-4x}$
\item $f(x) - 4 = \answer{-4x-1}$
\item  $f\left(x^2\right) = \answer{3-4x^2}$
\end{itemize}

\end{problem}  

\begin{problem}
$f(x) = 2 - x^2$

\begin{itemize}
\item $f(3) = \answer{-7}$
\item $f(-1) = \answer{1}$
\item $f\left(\frac{3}{2} \right) = \answer{-\frac{1}{4}}$
\item  $f(4x) = \answer{2-16x^2}$
\item $4f(x) = \answer{8-4x^2}$
\item $f(-x) = \answer{2-x^2}$
\item  $f(x-4) = \answer{-x^2+8x-14}$
\item $f(x) - 4 = \answer{-x^2-2}$
\item  $f\left(x^2\right) = \answer{2-x^4}$
\end{itemize}

\end{problem}  


\begin{problem}
$f(x) = x^2 - 3x + 2$,


\begin{itemize}
\item $f(3) = \answer{2}$
\item $f(-1) = \answer{6}$
\item $f\left(\frac{3}{2} \right) = \answer{-\frac{1}{4}}$
\item  $f(4x) = \answer{16x^2-12x+2}$
\item $4f(x) = \answer{4x^2-12x+8}$
\item $f(-x) = \answer{x^2+3x+2}$
\item  $f(x-4) = \answer{x^2-11x+30}$
\item $f(x) - 4 = \answer{x^2-3x-2}$
\item  $f\left(x^2\right) = \answer{x^4-3x^2+2}$
\end{itemize}  
\end{problem}

\begin{problem}
$f(x) = 6$

\begin{itemize}
\item $f(3) = \answer{6}$
\item $f(-1) =\answer{6}$
\item $f\left(\frac{3}{2} \right) = \answer{6}$
\item  $f(4x) = \answer{6}$
\item $4f(x) = \answer{24}$
\item $f(-x) = \answer{6}$
\item  $f(x-4) = \answer{6}$ 
\item $f(x) - 4 = \answer{2}$
\item  $f\left(x^2\right) = \answer{6}$
\end{itemize}

\end{problem}  

\begin{problem}\label{funcnotationbasiclast}
$f(x) = 0$

\begin{itemize}
\item $f(3) = \answer{0}$
\item $f(-1) =\answer{0}$
\item $f\left(\frac{3}{2} \right) = \answer{0}$
\item  $f(4x) = \answer{0}$
\item $4f(x) = \answer{0}$
\item $f(-x) = \answer{0}$
\item  $f(x-4) = \answer{0}$ 
\item $f(x) - 4 = \answer{-4}$
\item  $f\left(x^2\right) = \answer{0}$
\end{itemize}

\end{problem}  

\end{question}

\begin{question}
In Exercises \ref{secondfuncnotationbasicfirst} - \ref{secondfuncnotationbasiclast}, use the given function $f$ to find and simplify the following:

\begin{itemize}
\item  $f(2)$
\item  $f(-2)$
\item  $f(2a)$
\item  $2 f(a)$
\item $f(a+2)$
\item $f(a) + f(2)$
\item  $f \left( \frac{2}{a} \right)$
\item $\frac{f(a)}{2}$
\item  $f(a + h)$
\end{itemize}
    
\begin{problem}\label{secondfuncnotationbasicfirst}
$f(x) = 2x-5$

\begin{itemize}
\item  $f(2)=\answer{-1}$
\item  $f(-2)=\answer{-9}$
\item  $f(2a)=\answer{4a-5}$
\item  $2 f(a)=\answer{4a-10}$
\item $f(a+2)=\answer{2a-1}$
\item $f(a) + f(2)=\answer{2a-6}$
\item  $f \left( \frac{2}{a} \right)=\answer{\frac{4}{a} - 5}$
\item $\frac{f(a)}{2}=\answer{\frac{2a-5}{2}}$
\item  $f(a + h)=\answer{2a + 2h - 5}$
\end{itemize}

\end{problem}

\begin{problem}
$f(t) = 5-2t$

\begin{itemize}
\item  $f(2)=\answer{1}$
\item  $f(-2)=\answer{9}$
\item  $f(2a)=\answer{5-4a}$
\item  $2 f(a)=\answer{10-4a}$
\item $f(a+2)=\answer{1-2a}$
\item $f(a) + f(2)=\answer{6-2a}$
\item  $f \left( \frac{2}{a} \right)=\answer{5 - \frac{4}{a}}$
\item $\frac{f(a)}{2}=\answer{\frac{5-2a}{2}}$
\item  $f(a + h)=\answer{5-2a-2h}$
\end{itemize}

\end{problem}


\begin{problem}
$f(w) = 2w^2 - 1$

\begin{itemize}
\item  $f(2)=\answer{7}$
\item  $f(-2)=\answer{7}$
\item  $f(2a)=\answer{8a^2-1}$
\item  $2 f(a)=\answer{4a^2-2}$
\item $f(a+2)=\answer{2a^2+8a+7}$
\item $f(a) + f(2)=\answer{2a^2+6}$
\item  $f \left( \frac{2}{a} \right)=\answer{\frac{8}{a^2} - 1}$
\item $\frac{f(a)}{2}=\answer{\frac{2a^2-1}{2}}$
\item  $f(a + h)=\answer{2a^2+4ah+2h^2-1}$
\end{itemize}

\end{problem}


\begin{problem}
$f(q) = 3q^2+3q-2$


\begin{itemize}
\item  $f(2)=\answer{16}$
\item  $f(-2)=\answer{4}$
\item  $f(2a)=\answer{12a^2+6a-2}$
\item  $2 f(a)=\answer{6a^2+6a-4}$
\item $f(a+2)=\answer{3a^2+15a+16}$
\item $f(a) + f(2)=\answer{3a^2+3a+14}$
\item  $f \left( \frac{2}{a} \right)=\answer{\frac{12}{a^2} + \frac{6}{a} - 2}$
\item $\frac{f(a)}{2}=\answer{\frac{3a^2+3a-2}{2}}$
\item  $f(a + h)=\answer{3a^2 + 6ah + 3h^2+3a+3h-2}$
\end{itemize}

\end{problem}


\begin{problem}
$f(r) = 117$

\begin{itemize}
\item  $f(2)=\answer{117}$
\item  $f(-2)=\answer{117}$
\item  $f(2a)=\answer{117}$
\item  $2 f(a)=\answer{234}$
\item $f(a+2)=\answer{117}$
\item $f(a) + f(2)=\answer{117}$
\item  $f \left( \frac{2}{a} \right)=\answer{117}$
\item $\frac{f(a)}{2}=\answer{\frac{117}{2}}$
\item  $f(a + h)=\answer{117}$
\end{itemize}

\end{problem}
  

\begin{problem}\label{secondfuncnotationbasiclast}
$f(z) = \frac{z}{2}$

\begin{itemize}
\item  $f(2)=\answer{1}$
\item  $f(-2)=\answer{-1}$
\item  $f(2a)=\answer{a}$
\item  $2 f(a)=\answer{a}$
\item $f(a+2)=\answer{\frac{a+2}{2}}$
\item $f(a) + f(2)=\answer{\frac{a}{2}+ 1}$
\item  $f \left( \frac{2}{a} \right)=\answer{\frac{1}{a}}$
\item $\frac{f(a)}{2}=\answer{\frac{a}{4}}$
\item  $f(a + h)=\answer{\frac{a+h}{2}}$
\end{itemize}

\end{problem}
\end{question}

\begin{question}
In Exercises \ref{findzerofuncfirst} - \ref{findzerofunclast}, use the given function $f$ to find $f(0)$ and solve $f(x)=0$.

\begin{problem}\label{findzerofuncfirst}
$f(x) = 2x - 1$

$f(0) = \answer{-1}$

Solving $f(x)=0$ gives us 
\begin{selectAll}
    \choice{$x=-1$}
    \choice{$x=-\frac{1}{2}$}
    \choice{$x=0$}
    \choice[correct]{$x=\frac{1}{2}$}
    \choice{$x=1$}
\end{selectAll}
\end{problem}


\begin{problem}
$f(x) = 3 - \frac{2}{5} x$

$f(0) = \answer{3}$

Solving $f(x)=0$ gives us 
\begin{selectAll}
    \choice{$x=0$}
    \choice{$x=\frac{6}{5}$}
    \choice{$x=3$}
    \choice{$x=\frac{5}{13}$}
    \choice{$x=5$}
    \choice[correct]{$x=\frac{15}{2}$}
\end{selectAll}
\end{problem}


\begin{problem}
$f(x) = 2x^2 - 6$,

$f(0) = \answer{-6}$

Solving $f(x)=0$ gives us 
\begin{selectAll}
    \choice{$x=-6$}
    \choice{$x=-3$}
    \choice[correct]{$x=-\sqrt{3}$}
    \choice{$x=0$}
    \choice[correct]{$x=\sqrt{3}$}
    \choice{$x=3$}
    \choice{$x=6$}
\end{selectAll}
\end{problem}

\begin{problem}\label{findzerofunclast}
$f(x) = x^2 - x - 12$

$f(0) = \answer{-12}$

Solving $f(x)=0$ gives us 
\begin{selectAll}
    \choice{$x=-12$}
    \choice{$x=-4$}
    \choice[correct]{$x=-3$}
    \choice{$x=0$}
    \choice{$x=3$}
    \choice[correct]{$x=4$}
    \choice{$x=12$}
\end{selectAll}
\end{problem}

\end{question}

\begin{question}
In Exercises \ref{equfunctionfirst} - \ref{equfunctionlast}, determine whether or not the following equation represents $y$ as a function of $x$.

\begin{problem}\label{equfunctionfirst}
$y = x^{3} - x$

\begin{multipleChoice}
    \choice[correct]{Yes, $y$ is a function of $x$.}
    \choice{No, $y$ is a not a function of $x$.}
\end{multipleChoice}

\end{problem}

\begin{problem}
$y = \sqrt{x - 2}$

\begin{multipleChoice}
    \choice[correct]{Yes, $y$ is a function of $x$.}
    \choice{No, $y$ is a not a function of $x$.}
\end{multipleChoice}
\end{problem}

\begin{problem}
$x^{3}y = -4$ 

\begin{multipleChoice}
    \choice[correct]{Yes, $y$ is a function of $x$.}
    \begin{feedback}
        Yes, in fact, $y=\frac{-4}{x^3}$.
    \end{feedback}
    \choice{No, $y$ is a not a function of $x$.}
\end{multipleChoice}
\end{problem}

\begin{problem}
$x^{2} - y^{2} = 1$ 

\begin{multipleChoice}
    \choice{Yes, $y$ is a function of $x$.}
    \choice[correct]{No, $y$ is a not a function of $x$.}
    \begin{feedback}
        We cannot solve for y without using $\pm$ when we take the square root of both sides of our equation.
    \end{feedback}
\end{multipleChoice}
\end{problem}


\begin{problem}
$y = \frac{x}{x^{2} - 9}$ 

\begin{multipleChoice}
    \choice[correct]{Yes, $y$ is a function of $x$.}
    \choice{No, $y$ is a not a function of $x$.}
\end{multipleChoice}
\end{problem}


\begin{problem}
$x = -6$
\begin{multipleChoice}
    \choice{Yes, $y$ is a function of $x$.}
    \choice[correct]{No, $y$ is a not a function of $x$.}
\end{multipleChoice}
\end{problem}


\begin{problem}
$x = y^2 + 4$
\begin{multipleChoice}
    \choice{Yes, $y$ is a function of $x$.}
    \choice[correct]{No, $y$ is a not a function of $x$.}
    \begin{feedback}
        We cannot solve for y without using $\pm$ when we take the square root of both sides of our equation.
    \end{feedback}
\end{multipleChoice}
\end{problem}  

\begin{problem}
$y = x^2 + 4$
\begin{multipleChoice}
    \choice[correct]{Yes, $y$ is a function of $x$.}
    \choice{No, $y$ is a not a function of $x$.}
\end{multipleChoice}
\end{problem}

\begin{problem}
$x^2 + y^2 = 4$
\begin{multipleChoice}
    \choice{Yes, $y$ is a function of $x$.}
    \choice[correct]{No, $y$ is a not a function of $x$.}
    \begin{feedback}
        We cannot solve for y without using $\pm$ when we take the square root of both sides of our equation.
    \end{feedback}
\end{multipleChoice}
\end{problem}

\begin{problem}
$y = \sqrt{4-x^2}$
\begin{multipleChoice}
    \choice[correct]{Yes, $y$ is a function of $x$.}
    \choice{No, $y$ is a not a function of $x$.}
\end{multipleChoice}
\end{problem}

\begin{problem}
$x^2 - y^2 = 4$
\begin{multipleChoice}
    \choice{Yes, $y$ is a function of $x$.}
    \choice[correct]{No, $y$ is a not a function of $x$.}
    \begin{feedback}
        We cannot solve for y without using $\pm$ when we take the square root of both sides of our equation.
    \end{feedback}
\end{multipleChoice}
\end{problem}

\begin{problem}
$x^3 + y^3 = 4$
\begin{multipleChoice}
    \choice[correct]{Yes, $y$ is a function of $x$.}
    \begin{feedback}
        Yes, in fact, $y=\sqrt[3]{4-x^3}$.
    \end{feedback}
    \choice{No, $y$ is a not a function of $x$.}
\end{multipleChoice}
\end{problem}


\begin{problem}
$2x + 3y = 4$
\begin{multipleChoice}
    \choice[correct]{Yes, $y$ is a function of $x$.}
    \begin{feedback}
        Yes, in fact, $y=-\frac{2}{3}x+\frac{4}{3}$.
    \end{feedback}
    \choice{No, $y$ is a not a function of $x$.}
\end{multipleChoice}
\end{problem}

\begin{problem}
$2xy = 4$
\begin{multipleChoice}
    \choice[correct]{Yes, $y$ is a function of $x$.}
    \begin{feedback}
        Yes, in fact, $y=\frac{4}{2x}$.
    \end{feedback}
    \choice{No, $y$ is a not a function of $x$.}
\end{multipleChoice}
\end{problem}

\begin{problem}\label{equfunctionlast}
$x^2 = y^2$
\begin{multipleChoice}
    \choice{Yes, $y$ is a function of $x$.}
    \choice[correct]{No, $y$ is a not a function of $x$.}
    \begin{feedback}
        We cannot solve for y without using $\pm$ when we take the square root of both sides of our equation.
    \end{feedback}
\end{multipleChoice}
\end{problem}

\end{question}

\begin{question}
Exercises \ref{setfunctionfirst} - \ref{setfunctionlast} give a set of points in the $xy$-plane. Determine if $y$ is a function of $x$. If so, state the
domain and range.

\begin{problem}\label{setfunctionfirst}
\{$(-3, 9)$, $\;(-2, 4)$, $\;(-1, 1)$, $\;(0, 0)$, $\;(1, 1)$, $\;(2, 4)$, $\;(3, 9)\}$ 
\end{problem}

\begin{problem}
$\left\{ (-3,0), (1,6), (2,-3), (4,2), (-5,6), (4,-9), (6,2) \right\}$
\end{problem}    


\begin{problem}
$\left\{ (-3,0), (-7,6), (5,5), (6,4), (4,9), (3,0) \right\}$
\end{problem}    

\begin{problem}
$\left\{ (1,2), (4,4), (9,6), (16,8), (25,10), (36, 12), \ldots \right\}$
\end{problem}    

\begin{problem}
\{($x, y) \, | \, x$ is an odd integer, and $y$ is an even integer\}
\end{problem}   

\begin{problem}
\{$(x, 1) \, | \, x$ is an irrational number\}
\end{problem}   

\begin{problem}
$\{ (1,0), (2,1), (4,2), (8,3), (16,4), (32, 5), \ldots \}$
\end{problem}   

\begin{problem}
$\{ \ldots (-3,9), (-2,4), (-1,1), (0,0), (1,1), (2,4), (3,9), \ldots \}$
\end{problem}   

\begin{problem}
$\{ (-2, y) \, | \, -3 < y < 4\}$
\end{problem}   

\begin{problem}
$\{ (x,3) \, | \,  -2 \leq x < 4\}$
\end{problem}    


\begin{problem}
$\{ \left(x,x^2\right) \, | \, \text{$x$ is a real number} \}$
\end{problem}    

\begin{problem}\label{setfunctionlast}
$\{ \left(x^2,x\right) \, | \, \text{$x$ is a real number} \}$
\end{problem}     
    
\end{question}

\begin{problem}\label{HLTExercise} The Vertical Line Test is a quick way to determine from a graph if the vertical axis variable is a function of the horizontal axis variable. If we are given a graph and asked to determine if the horizontal axis variable is a function of the vertical axis variable, we can use horizontal lines instead of vertical lines to check.  Using Theorem \ref{VLT} as a guide,  formulate a `Horizontal Line Test.'  (We'll refer back to this exercise in Section \ref{InverseFunctions}.)
\end{problem}

%\setcounter{HW}{\value{enumi}}


 


In Exercises \ref{graphfunctionfirstxy} - \ref{graphfunctionlastxy}, determine whether or not the graph suggests $y$ is a function of $x$.  For the ones which do, state the domain and range. 


\begin{problem}\label{graphfunctionfirstxy}
    Determine whether or not the graph suggests $y$ is a function of $x$.  If it does, state the domain and range.

    (NEED GRAPH HERE)
\end{problem}

% \begin{mfpic}[17]{-5}{2}{-2}{5}
% \point[4pt]{(-4, -1), (-3, 0), (-2, 1), (-1, 2), (0, 3), (1, 4)}
% \axes
% \tlabel[cc](2,-0.5){\scriptsize $x$}
% \tlabel[cc](0.5,4.75){\scriptsize $y$}
% \xmarks{-4,-3,-2,-1,1}
% \ymarks{-1,1,2,3,4}
% \tlpointsep{4pt}
% \axislabels {x}{{\tiny $-4 \hspace{8pt}$} -4, {\tiny $-3 \hspace{8pt}$} -3, {\tiny $-2 \hspace{8pt}$} -2, {\tiny $-1 \hspace{8pt}$} -1, {\tiny $1$} 1}
% \axislabels {y}{{\tiny $-1$} -1, {\tiny $1$} 1, {\tiny $2$} 2, {\tiny $3$} 3, {\tiny $4$} 4}
% \end{mfpic}

%  
%  

\begin{problem}\label{graphfunctionfirstxy2}
    Determine whether or not the graph suggests $y$ is a function of $x$.  If it does, state the domain and range.

    (NEED GRAPH HERE)
\end{problem} 

% \begin{mfpic}[15]{-5}{2}{-2}{5}
% \point[4pt]{(-4, -1), (-3, 0), (-3, 1), (-2, 1), (-1, 2), (0, 3), (1, 4)}
% \axes
% \tlabel[cc](2,-0.5){\scriptsize $x$}
% \tlabel[cc](0.5,4.75){\scriptsize $y$}
% \xmarks{-4,-3,-2,-1,1}
% \ymarks{-1,1,2,3,4}
% \tlpointsep{4pt}
% \axislabels {x}{{\tiny $-4 \hspace{6pt}$} -4, {\tiny $-3 \hspace{6pt}$} -3, {\tiny $-2 \hspace{6pt}$} -2, {\tiny $-1 \hspace{6pt}$} -1, {\tiny $1$} 1}
% \axislabels {y}{{\tiny $-1$} -1, {\tiny $1$} 1, {\tiny $2$} 2, {\tiny $3$} 3, {\tiny $4$} 4}
% \end{mfpic}


%\setcounter{HW}{\value{enumi}}

%\setcounter{enumi}{\value{HW}}

\begin{problem}\label{graphfunctionfirstxy3}
    Determine whether or not the graph suggests $y$ is a function of $x$.  If it does, state the domain and range.

    (NEED GRAPH HERE)
\end{problem}

% \begin{mfpic}[15]{-3}{3}{-1}{6}
% \axes
% \tlabel[cc](3,-0.5){\scriptsize $x$}
% \tlabel[cc](0.5,5.75){\scriptsize $y$}
% \xmarks{-2,-1,1,2}
% \ymarks{1,2,3,4,5}
% \tlpointsep{4pt}
% \axislabels {x}{{\tiny $-2 \hspace{8pt}$} -2, {\tiny $-1 \hspace{8pt}$} -1, {\tiny $1$} 1, {\tiny $2$} 2}
% \axislabels {y}{{\tiny $1$} 1, {\tiny $2$} 2, {\tiny $3$} 3, {\tiny $4$} 4, {\tiny $5$} 5}
% \penwd{1.25pt}
% \arrow \reverse \arrow \function{-2.1, 2.1, 0.1}{x**2+1}
% \end{mfpic}

 
\begin{problem}\label{graphfunctionlastxy}
    Determine whether or not the graph suggests $y$ is a function of $x$.  If it does, state the domain and range.

    (NEED GRAPH HERE)
\end{problem} 

% \begin{mfpic}[15]{-4}{4}{-4}{4}
% \axes
% \tlabel[cc](4,-0.5){\scriptsize $x$}
% \tlabel[cc](0.5,3.75){\scriptsize $y$}
% \xmarks{-3,-2,-1,1,2,3}
% \ymarks{-3,-2,-1,1,2,3}
% \tlpointsep{4pt}
% \axislabels {x}{{\tiny $-3 \hspace{8pt}$} -3, {\tiny $-2 \hspace{8pt}$} -2, {\tiny $-1 \hspace{8pt}$} -1, {\tiny $1$} 1, {\tiny $2$} 2, {\tiny $3$} 3}
% \axislabels {y}{{\tiny $-3$} -3, {\tiny $-2$} -2, {\tiny $-1$} -1, {\tiny $1$} 1, {\tiny $2$} 2, {\tiny $3$} 3}
% \penwd{1.25pt}
% \arrow \reverse \arrow \parafcn{-2,2,0.1}{(cosh(t),sinh(t))}
% \end{mfpic}


%\setcounter{HW}{\value{enumi}}

%\setcounter{enumi}{\value{HW}}

\begin{problem}
    Determine whether or not the graph suggests $x$ is a function of $y$.  (Feel free to use Exercise \ref{HLTExercise}.)  If it does, state the domain and range.

    (NEED 4 GRAPHS HERE)
\end{problem}

  

%\setcounter{HW}{\value{enumi}}



%\setcounter{enumi}{\value{HW}}

  

\begin{problem}\label{graphfunctionfirstvw}
    Determine whether or not the graph suggests $w$ is a function of $v$  If it does, state the domain and range.

    (NEED GRAPH HERE)
\end{problem} 
% \begin{mfpic}[15]{-1}{10}{-1}{4}
% \axes
% \tlabel[cc](10,-0.5){\scriptsize $v$}
% \tlabel[cc](0.5,3.75){\scriptsize $w$}
% \xmarks{1,2,3,4,5,6,7,8,9}
% \ymarks{1,2,3}
% \tlpointsep{4pt}
% \axislabels {x}{{\tiny $1$} 1, {\tiny $2$} 2, {\tiny $3$} 3, {\tiny $4$} 4, {\tiny $5$} 5, {\tiny $6$} 6, {\tiny $7$} 7, {\tiny $8$} 8, {\tiny $9$} 9}
% \axislabels {y}{{\tiny $1$} 1, {\tiny $2$} 2, {\tiny $3$} 3}
% \penwd{1.25pt}
% \arrow \function{2, 10, 0.1}{sqrt(x - 2)}
% \point[4pt]{(2,0)}
% \end{mfpic}


\begin{problem}\label{graphfunctionfirstvw2}
    Determine whether or not the graph suggests $w$ is a function of $v$  If it does, state the domain and range.

    (NEED GRAPH HERE)
\end{problem} 

% \begin{mfpic}[15]{-5}{5}{-1}{5}
% \axes
% \tlabel[cc](5,-0.5){\scriptsize $v$}
% \tlabel[cc](0.5,4.75){\scriptsize $w$}
% \xmarks{-4,-3,-2,-1,1,2,3,4}
% \ymarks{1,2,3,4}
% \tlpointsep{4pt}
% \axislabels {x}{{\tiny $-4 \hspace{8pt}$} -4, {\tiny $-3 \hspace{8pt}$} -3, {\tiny $-2 \hspace{8pt}$} -2, {\tiny $-1 \hspace{8pt}$} -1, {\tiny $1$} 1, {\tiny $2$} 2, {\tiny $3$} 3, {\tiny $4$} 4}
% \axislabels {y}{{\tiny $1$} 1, {\tiny $2$} 2, {\tiny $3$} 3, {\tiny $4$} 4}
% \penwd{1.25pt}
% \arrow \reverse \arrow \function{-5, 5, 0.1}{4/(x**2 + 1)}
% \end{mfpic}


%\setcounter{HW}{\value{enumi}}

%\setcounter{enumi}{\value{HW}}

\begin{problem}\label{graphfunctionfirstvw3}
    Determine whether or not the graph suggests $w$ is a function of $v$  If it does, state the domain and range.

    (NEED GRAPH HERE)
\end{problem}    

% \begin{mfpic}[15]{-4.5}{5.5}{-4}{3}
% \axes
% %\fillcolor[gray]{.7}
% %\gfill \rect{(-3.97, -2.97), (4.97, 1.97)}
% %\dashed \polyline{(-4, 2), (5, 2)}
% %\dashed \polyline{(5, 2), (5, -3)}
% %\dashed \polyline{(5, -3), (-4, -3)}
% \tlabel[cc](5.5,-0.5){\scriptsize $v$}
% \tlabel[cc](0.5,2.75){\scriptsize $w$}
% \xmarks{-4,-3,-2,-1,1,2,3,4,5}
% \ymarks{-3,-2,-1,1,2}
% \tlpointsep{4pt}
% \axislabels {x}{{\tiny $-4 \hspace{8pt}$} -4,  {\tiny $-2 \hspace{8pt}$} -2, {\tiny $-1 \hspace{8pt}$} -1, {\tiny $1$} 1, {\tiny $2$} 2, {\tiny $4$} 4, {\tiny $5$} 5}
% \axislabels {y}{ {\tiny $-2$} -2, {\tiny $-1$} -1, {\tiny $1$} 1}
% \penwd{1.25pt}
% \polyline{(-3, -3), (-3, 2), (3,2), (3,-3), (-3,-3)}
% \end{mfpic}



\begin{problem}\label{graphfunctionlastvw}
    Determine whether or not the graph suggests $w$ is a function of $v$  If it does, state the domain and range.

    (NEED GRAPH HERE)
\end{problem} 

% \begin{mfpic}[15]{-6}{4}{-2.5}{4.5}
% \axes
% \tlabel[cc](4,-0.5){\scriptsize $v$}
% \tlabel[cc](0.5,4.75){\scriptsize $w$}
% \xmarks{-5 step 1 until 3}
% \ymarks{-2 step 1 until 4}
% \tlpointsep{4pt}
% \axislabels {x}{{\tiny $-5 \hspace{8pt}$} -5, {\tiny $-4 \hspace{8pt}$} -4, {\tiny $-3 \hspace{8pt}$} -3, {\tiny $-2 \hspace{8pt}$} -2, {\tiny $-1 \hspace{8pt}$} -1, {\tiny $1$} 1, {\tiny $2$} 2, {\tiny $3$} 3}
% \axislabels {y}{{\tiny $-2$} -2, {\tiny $-1$} -1, {\tiny $1$} 1, {\tiny $2$} 2, {\tiny $3$} 3, {\tiny $4$} 4}
% \penwd{1.25pt}
% \function{-5,-1,0.1}{-5 - 6*x - x**2}
% \function{-1,3,0.1}{x/4 - 7/4}
% \point[4pt]{(-5, 0), (-1, 0)}
% \pointfillfalse
% \point[4pt]{(-3,4), (-1,-2), (3,-1)}
% \end{mfpic}


%\setcounter{HW}{\value{enumi}}

%\setcounter{enumi}{\value{HW}}

\begin{problem}
    Determine whether or not the graph suggests $v$ is a function of $w$.  (Feel free to use Exercise \ref{HLTExercise}.)  If it does, state the domain and range.

    (NEED 4 GRAPHS HERE)
\end{problem}

%\setcounter{HW}{\value{enumi}}


%In Exercises \ref{graphfunctionfirsttT} - \ref{graphfunctionlasttT}, determine whether or not the graph suggests $T$ is a function of $t$.   For the ones which do, state the domain and range. 



%\setcounter{enumi}{\value{HW}}

\begin{problem}\label{graphfunctionfirsttT}
    Determine whether or not the graph suggests $T$ is a function of $t$  If it does, state the domain and range.

    (NEED GRAPH HERE)
\end{problem} 

% \begin{mfpic}[8]{-4}{4}{-6}{10}

% \axes
% \tlabel[cc](4,-0.5){\scriptsize $t$}
% \tlabel[cc](0.5,10){\scriptsize $T$}
% \xmarks{-3,-2,-1,1,2,3}
% \ymarks{-5,-4,-3,-2,-1,1,2,3,4,5,6,7,8,9}
% \tlpointsep{4pt}
% \axislabels {x}{{\tiny $-3 \hspace{6pt}$} -3,{\tiny $-2 \hspace{6pt}$} -2, {\tiny $-1 \hspace{6pt}$} -1, {\tiny $1$} 1, {\tiny $2$} 2, {\tiny $3$} 3}
% \axislabels {y}{{\tiny $-5$} -5, {\tiny $-4$} -4, {\tiny $-3$} -3, {\tiny $-2$} -2, {\tiny $-1$} -1, {\tiny $1$} 1, {\tiny $2$} 2, {\tiny $3$} 3, {\tiny $4$} 4, {\tiny $5$} 5, {\tiny $6$} 6, {\tiny $7$} 7, {\tiny $8$} 8, {\tiny $9$} 9}
% \penwd{1.25pt}
% \arrow \function{-2,4.5,0.1}{x**2 - 2*x - 2}
% \point[4pt]{(-2,6), (1,-3) }
% \end{mfpic}


\begin{problem}\label{graphfunctionfirsttT2}
    Determine whether or not the graph suggests $T$ is a function of $t$  If it does, state the domain and range.

    (NEED GRAPH HERE)
\end{problem} 

% \begin{mfpic}[10]{-6}{6}{-6}{6}
% \axes
% \tlabel[cc](6,-0.5){\scriptsize $t$}
% \tlabel[cc](0.5,6){\scriptsize $T$}
% \xmarks{-5,-4,-3,-2,-1,1,2,3,4,5}
% \ymarks{-5,-4,-3,-2,-1,1,2,3,4,5}
% \tlpointsep{4pt}
% \axislabels {x}{{\tiny $-5 \hspace{6pt}$} -5,{\tiny $-4 \hspace{6pt}$} -4,{\tiny $-3 \hspace{6pt}$} -3,{\tiny $-2 \hspace{6pt}$} -2, {\tiny $-1 \hspace{6pt}$} -1, {\tiny $1$} 1, {\tiny $2$} 2, {\tiny $3$} 3, {\tiny $4$} 4, {\tiny $5$} 5}
% \penwd{1.25pt}
% \plrfcn{0,180,5}{5*sind 3t}
% \end{mfpic} 


\begin{problem}\label{graphfunctionfirsttT3}
    Determine whether or not the graph suggests $T$ is a function of $t$  If it does, state the domain and range.

    (NEED GRAPH HERE)
\end{problem} 

% \begin{mfpic}[10]{-6}{6}{-6}{6}
% \axes
% \tlabel[cc](6,-0.5){\scriptsize $t$}
% \tlabel[cc](0.5,6){\scriptsize $T$}
% \xmarks{-5,-4,-3,-2,-1,1,2,3,4,5}
% \ymarks{-5,-4,-3,-2,-1,1,2,3,4,5}
% \tlpointsep{4pt}
% \axislabels {x}{{\tiny $-5 \hspace{6pt}$} -5,{\tiny $-4 \hspace{6pt}$} -4,{\tiny $-3 \hspace{6pt}$} -3,{\tiny $-2 \hspace{6pt}$} -2, {\tiny $-1 \hspace{6pt}$} -1, {\tiny $1$} 1, {\tiny $2$} 2, {\tiny $3$} 3, {\tiny $4$} 4, {\tiny $5$} 5}
% \axislabels {y}{{\tiny $-5$} -5,{\tiny $-4$} -4,{\tiny $-3$} -3, {\tiny $-2$} -2, {\tiny $-1$} -1, {\tiny $1$} 1, {\tiny $2$} 2, {\tiny $3$} 3, {\tiny $4$} 4, {\tiny $5$} 5}
% \penwd{1.25pt}
% \function{-5,4,0.1}{0.0502*(x**3) - 0.0344*(x**2) - 0.2010*x + 2.138}
% \pointfillfalse
% \point[4pt]{(-5,-4),(4,4)}

% \end{mfpic} 



\begin{problem}\label{graphfunctionlasttT}
    Determine whether or not the graph suggests $T$ is a function of $t$  If it does, state the domain and range.

    (NEED GRAPH HERE)
\end{problem} 

% \begin{mfpic}[10]{-2}{7}{-6}{6}
% \axes
% \tlabel[cc](7,-0.5){\scriptsize $t$}
% \tlabel[cc](0.5,6){\scriptsize $T$}
% \xmarks{-1,1,2,3,4,5,6}
% \ymarks{-5,-4,-3,-2,-1,1,2,3,4,5}
% \tlpointsep{4pt}
% \axislabels {x}{{\tiny $-1 \hspace{6pt}$} -1, {\tiny $1$} 1, {\tiny $2$} 2, {\tiny $3$} 3, {\tiny $4$} 4, {\tiny $5$} 5, {\tiny $6$} 6}
% \axislabels {y}{{\tiny $-5$} -5,{\tiny $-4$} -4,{\tiny $-3$} -3, {\tiny $-2$} -2, {\tiny $-1$} -1, {\tiny $1$} 1, {\tiny $2$} 2, {\tiny $3$} 3, {\tiny $4$} 4, {\tiny $5$} 5}
% \penwd{1.25pt}
% \polyline{(0,-1), (3,-4)}
% \polyline{(3,1), (4,4), (6,0)}
% \point[4pt]{(0,-1), (4,4), (6,0)}
% \pointfillfalse
% \point[4pt]{(3,-4), (3,1)}
% \end{mfpic} 

%\setcounter{HW}{\value{enumi}}



%\setcounter{enumi}{\value{HW}}

\begin{problem}
    Determine whether or not the graph suggests $t$ is a function of $T$.  (Feel free to use Exercise \ref{HLTExercise}.)  If it does, state the domain and range.

    (NEED 4 GRAPHS HERE)
\end{problem}


%\setcounter{HW}{\value{enumi}}



In Exercises \ref{graphfunctionfirstHs} - \ref{graphfunctionlastHs}, determine whether or not the graph suggests $H$ is a function of $s$.  For the ones which do, state the domain and range. 


\begin{enumerate}
%\setcounter{enumi}{\value{HW}}

\item  $~$ \label{graphfunctionfirstHs}

% \begin{mfpic}[15]{-3}{3}{-1}{5}
% \axes
% \tlabel[cc](3,-0.5){\scriptsize $s$}
% \tlabel[cc](0.5,5){\scriptsize $H$}
% \xmarks{-2,-1,1,2}
% \ymarks{1,2,3,4}
% \tlpointsep{4pt}
% \axislabels {x}{{\tiny $-2 \hspace{6pt}$} -2, {\tiny $-1 \hspace{6pt}$} -1, {\tiny $1$} 1, {\tiny $2$} 2}
% \axislabels {y}{{\tiny $1$} 1, {\tiny $2$} 2, {\tiny $3$} 3, {\tiny $4$} 4}
% \penwd{1.25pt}
% \arrow \reverse \arrow \function{-2.25,2.25,0.1}{4-(x**2)}
% \end{mfpic} 


\item  $~$  \label{graphfunctionfirstHs2}


% \begin{mfpic}[15]{-3}{3}{-1}{5}
% \axes
% \tlabel[cc](3,-0.5){\scriptsize $s$}
% \tlabel[cc](0.5,5){\scriptsize $H$}
% \xmarks{-2,-1,1,2}
% \ymarks{1,2,3,4}
% \tlpointsep{4pt}
% \axislabels {x}{{\tiny $-2 \hspace{6pt}$} -2, {\tiny $-1 \hspace{6pt}$} -1, {\tiny $1$} 1, {\tiny $2$} 2}
% \axislabels {y}{{\tiny $1$} 1, {\tiny $2$} 2, {\tiny $3$} 3, {\tiny $4$} 4}
% \penwd{1.25pt}
% \arrow \reverse \arrow \polyline{(-2,-1), (1,4), (2,-1)}
% \end{mfpic} 

%\setcounter{HW}{\value{enumi}}
\end{enumerate}


\begin{enumerate}
%\setcounter{enumi}{\value{HW}}

\item  $~$  \label{graphfunctionfirstHs3}

% \begin{mfpic}[15]{-3}{3}{-1}{5}
% \axes
% \tlabel[cc](3,-0.5){\scriptsize $s$}
% \tlabel[cc](0.5,5){\scriptsize $H$}
% \xmarks{-2,-1,1,2}
% \ymarks{1,2,3,4}
% \tlpointsep{4pt}
% \axislabels {x}{{\tiny $-2 \hspace{6pt}$} -2, {\tiny $-1 \hspace{6pt}$} -1, {\tiny $1$} 1, {\tiny $2$} 2}
% \axislabels {y}{{\tiny $1$} 1, {\tiny $2$} 2, {\tiny $3$} 3, {\tiny $4$} 4}
% \penwd{1.25pt}
% \arrow \function{-2, 1.8, 0.1}{3-2*sqrt(x+2)}
% \point[4pt]{(-2,3)}
% \end{mfpic} 

 

\item  $~$ \label{graphfunctionlastHs}


% \begin{mfpic}[15]{-3}{3}{-1}{5}
% \axes
% \tlabel[cc](3,-0.5){\scriptsize $s$}
% \tlabel[cc](0.5,5){\scriptsize $H$}
% \xmarks{-2,-1,1,2}
% \ymarks{1,2,3,4}
% \tlpointsep{4pt}
% \axislabels {x}{{\tiny $-2 \hspace{6pt}$} -2, {\tiny $-1 \hspace{6pt}$} -1, {\tiny $1$} 1, {\tiny $2$} 2}
% \axislabels {y}{{\tiny $1$} 1, {\tiny $2$} 2, {\tiny $3$} 3, {\tiny $4$} 4}
% \penwd{1.25pt}
% \arrow \reverse \arrow \function{-2.15, 1.75, 0.1}{x*(x-1)*(x+2)}
% \end{mfpic} 

%\setcounter{HW}{\value{enumi}}
\end{enumerate}


\begin{enumerate}
%\setcounter{enumi}{\value{HW}}

\item  Determine which, if any, of the graphs in numbers \ref{graphfunctionfirstHs} - \ref{graphfunctionlastHs} represent $s$ as a function of $H$.  For the ones which do, state the domain and range.   (Feel free to use Exercise \ref{HLTExercise}.)

%\setcounter{HW}{\value{enumi}}
\end{enumerate}


In Exercises \ref{graphfunctionfirstut} - \ref{graphfunctionlastut}, determine whether or not the graph suggests $u$ is a function of $t$. For the ones which do, state the domain and range. 



\begin{enumerate}
%\setcounter{enumi}{\value{HW}}

\item  $~$   \label{graphfunctionfirstut}

% \begin{mfpic}[15]{-3}{3}{-3}{3}
% \axes
% \tlabel[cc](3,-0.5){\scriptsize $t$}
% \tlabel[cc](0.5,3){\scriptsize $u$}
% \xmarks{-2,-1,1,2}
% \ymarks{-2,-1,1,2}
% \tlpointsep{4pt}
% \axislabels {x}{{\tiny $-2 \hspace{6pt}$} -2, {\tiny $-1 \hspace{6pt}$} -1, {\tiny $1$} 1, {\tiny $2$} 2}
% \axislabels {y}{{\tiny $1$} 1, {\tiny $2$} 2, {\tiny $-2$} -2, {\tiny $-1$} -1}
% \penwd{1.25pt}
% \arrow \polyline{(0,1), (-2,-2)}
% \arrow \polyline{(1,2), (3,2)}
% \point[4pt]{(0,1)}
% \pointfillfalse
% \point[4pt]{(1,2)}
% \end{mfpic} 

%  
%  

\item  $~$  \label{graphfunctionfirstut2}


% \begin{mfpic}[15]{-4}{4}{-3}{3}
% \axes
% \tlabel[cc](4,-0.5){\scriptsize $t$}
% \tlabel[cc](0.5,3){\scriptsize $u$}
% \xmarks{-3,-2,-1,1,2,3}
% \ymarks{-2,-1,1,2}
% \tlpointsep{4pt}
% \axislabels {x}{{\tiny $-3 \hspace{6pt}$} -3,{\tiny $-2 \hspace{6pt}$} -2, {\tiny $-1 \hspace{6pt}$} -1, {\tiny $1$} 1, {\tiny $2$} 2, {\tiny $3$} 3}
% \axislabels {y}{{\tiny $1$} 1, {\tiny $2$} 2, {\tiny $-2$} -2, {\tiny $-1$} -1}
% \penwd{1.25pt}
% \function{-3,3,0.1}{2*sin(1.05*x)}
% \point[4pt]{(-3,0)}
% \point[4pt]{(3,0)}
% \end{mfpic} 

%\setcounter{HW}{\value{enumi}}
\end{enumerate}



\begin{enumerate}
%\setcounter{enumi}{\value{HW}}

\item  $~$  \label{graphfunctionfirstut3}

% \begin{mfpic}[15]{-3}{3}{-3}{3}
% \axes
% \tlabel[cc](3,-0.5){\scriptsize $t$}
% \tlabel[cc](0.5,3){\scriptsize $u$}
% \xmarks{-2,-1,1,2}
% \ymarks{-2,-1,1,2}
% \tlpointsep{4pt}
% \axislabels {x}{{\tiny $-2 \hspace{6pt}$} -2, {\tiny $-1 \hspace{6pt}$} -1, {\tiny $1$} 1, {\tiny $2$} 2}
% \axislabels {y}{{\tiny $1$} 1, {\tiny $2$} 2, {\tiny $-2$} -2, {\tiny $-1$} -1}
% \penwd{1.25pt}
% \arrow \reverse \arrow \polyline{(2,-3), (2,3)}
% \end{mfpic} 

%  
%  

\item  $~$ \label{graphfunctionlastut}


% \begin{mfpic}[15]{-3}{3}{-3}{3}
% \axes
% \tlabel[cc](3,-0.5){\scriptsize $t$}
% \tlabel[cc](0.5,3){\scriptsize $u$}
% \xmarks{-2,-1,1,2}
% \ymarks{-2,-1,1,2}
% \tlpointsep{4pt}
% \axislabels {x}{{\tiny $-2 \hspace{6pt}$} -2, {\tiny $-1 \hspace{6pt}$} -1, {\tiny $1$} 1, {\tiny $2$} 2}
% \axislabels {y}{{\tiny $1$} 1, {\tiny $2$} 2, {\tiny $-2$} -2, {\tiny $-1$} -1}
% \penwd{1.25pt}
% \arrow \reverse \arrow \polyline{(-3,2), (3,2)}
% \end{mfpic} 

%\setcounter{HW}{\value{enumi}}
\end{enumerate}


\begin{enumerate}
%\setcounter{enumi}{\value{HW}}

\item  Determine which, if any, of the graphs in numbers \ref{graphfunctionfirstut} - \ref{graphfunctionlastut} represent $t$ as a function of $u$.  For the ones which do, state the domain and range.   (Feel free to use Exercise \ref{HLTExercise}.)

%\setcounter{HW}{\value{enumi}}
\end{enumerate}



In Exercises \ref{functionvaluesfromgraphfirst} - \ref{functionvaluesfromgraphlast}, use the graphs of $f$ and $g$ below to find the indicated values.



% \begin{mfpic}[15]{-6}{6}{-6}{6}

% \axes
% \tlabel[cc](6,-0.5){\scriptsize $x$}
% \tlabel[cc](0.5,6){\scriptsize $y$}
% \xmarks{-5,-4,-3,-2,-1,1,2,3,4,5}
% \ymarks{-5,-4,-3,-2,-1,1,2,3,4,5}
% \tlpointsep{5pt}
% \scriptsize
% \axislabels {x}{{$-5 \hspace{7pt}$} -5,{$-4 \hspace{7pt}$} -4,{$-3 \hspace{7pt}$} -3,{$-2 \hspace{7pt}$} -2, {$-1 \hspace{7pt}$} -1, {$1$} 1, {$2$} 2, {$3$} 3, {$4$} 4, {$5$} 5}
% \axislabels {y}{{$-5$} -5,{$-4$} -4,{$-3$} -3,{$-2$} -2,{$-1$} -1, {$1$} 1, {$2$} 2, {$3$} 3, {$4$} 4, {$5$} 5}
% \normalsize
% \tcaption{The graph of $y = f(x)$.}
% \penwd{1.25pt}
% \point[4pt]{(-5, -5), (-4, 0), (-3, 4), (-2,2), (-1,0), (0,-1), (1,0), (2,3), (3,1)}
% \polyline{(-5,-5), (-4,0), (-3,4), (-2,2), (-1,0)}
% \function{-1, 1, 0.1}{(x**2)-1}
% \polyline{(1,0), (2,3), (3,1)}
% \end{mfpic}



% \begin{mfpic}[15]{-5}{5}{-6}{6}

% \axes
% \tlabel[cc](5,-0.5){\scriptsize $t$}
% \tlabel[cc](0.5,6){\scriptsize $s$}
% \xmarks{-3,-2,-1,1,2,3,4}
% \ymarks{-5,-4,-3,-2,-1,1,2,3,4,5}
% \tlpointsep{5pt}
% \scriptsize
% \axislabels {x}{{$-4 \hspace{7pt}$} -4, {$-3 \hspace{7pt}$} -3,{$-2 \hspace{7pt}$} -2, {$-1 \hspace{7pt}$} -1, {$1$} 1, {$2$} 2, {$3$} 3, {$4$} 4}
% \axislabels {y}{{$-5$} -5,{$-4$} -4,{$-3$} -3,{$-2$} -2,{$-1$} -1, {$1$} 1, {$2$} 2, {$3$} 3, {$4$} 4, {$5$} 5}
% \normalsize
% \tcaption{The graph of $s = g(t)$.}
% \penwd{1.25pt}
% \function{-4, 4, 0.1}{5*sin(x*3.14159/4)}
% \point[4pt]{ (-2, -5), (0, 0), (4,0), (2,3), (-4,0)}
% \pointfillfalse
% \point[4pt]{(2,5)}
% \end{mfpic}


\begin{enumerate}

%\setcounter{enumi}{\value{HW}}

\item \label{functionvaluesfromgraphfirst} $f(-2)$

\item $g(-2)$

\item $f(2)$

\item  $g(2)$

%\setcounter{HW}{\value{enumi}}

\end{enumerate}




\begin{enumerate}

%\setcounter{enumi}{\value{HW}}

\item $f(0)$

\item $g(0)$

\item  Solve $f(x) = 0$.

\item  Solve $g(t) = 0$. 

%\setcounter{HW}{\value{enumi}}

\end{enumerate}




\begin{enumerate}

%\setcounter{enumi}{\value{HW}}

\item  State the domain and range of $f$.

\item  State the domain and range of $g$ . \label{functionvaluesfromgraphlast}

%\setcounter{HW}{\value{enumi}}

\end{enumerate}


 

In Exercises \ref{sketchgraphfirst} - \ref{sketchgraphlast}, graph each function by making a table, plotting points, and using a graphing utility (if needed.)  Use the independent variable as the horizontal axis label and the default `$y$' label for the vertical axis label.  State the domain and range of each function. 


\begin{enumerate}
%\setcounter{enumi}{\value{HW}}
\item $f(x) = 2-x$ \label{sketchgraphfirst}
\item $g(t) = \frac{t - 2}{3}$
\item $h(s) = s^2 + 1$

%\setcounter{HW}{\value{enumi}}
\end{enumerate}


\begin{enumerate}
%\setcounter{enumi}{\value{HW}}

\item $f(x) = 4-x^2$
\item $g(t) = 2$
\item $h(s) = s^3$

%\setcounter{HW}{\value{enumi}}
\end{enumerate}


\begin{enumerate}
%\setcounter{enumi}{\value{HW}}

\item $f(x) = x(x-1)(x+2)$
\item $g(t) = \sqrt{t-2}$
\item $h(s) = \sqrt{5 - s}$

%\setcounter{HW}{\value{enumi}}
\end{enumerate}


\begin{enumerate}
%\setcounter{enumi}{\value{HW}}

\item $f(x) = 3-2\sqrt{x+2}$
\item $g(t) = \sqrt[3]{t}$
\item $h(s) = \frac{1}{s^{2} + 1}$ \label{sketchgraphlast}

%\setcounter{HW}{\value{enumi}}
\end{enumerate}



\begin{enumerate}
%\setcounter{enumi}{\value{HW}}


\item Consider the function $f$ described below:

%\begin{center}

% \begin{mfpic}[19]{-5}{5}{-5}{6}
% \tlabel[cc](-3,5){$-1$}
% \tlabel[cc](-3,3){0}
% \tlabel[cc](-3,1){1}
% \tlabel[cc](-3,-1){2}
% \tlabel[cc](0,6){$f$}
% \tlabel[cc](3.5,5){$-3$}
% \tlabel[cc](3.5,3){0}
% \tlabel[cc](3.5,1){4}
% \arrow[l 5pt] \polyline{(-2.5, 5), (2.5, 3)}
% \arrow[l 5pt] \polyline{(-2.5, 3), (2.5, 5)}
% \arrow[l 5pt] \polyline{(-2.5, 1), (2.5, 3)}
% \arrow[l 5pt] \polyline{(-2.5, -1), (2.5, 1)}
% \end{mfpic}

% \end{center}

\begin{enumerate}

\item  State  the domain and range of $f$.

\item Find $f(0)$ and solve $f(x) = 0$.

\item  Write $f$ as a set of ordered pairs.

\item  Graph $f$.

\end{enumerate}


\item  Let $g = \{ (-1,4), (0,2), (2, 3), (3,4)  \}$

\begin{enumerate}

\item  State the domain and range of $g$.

\item  Create a mapping diagram for $g$.

\item  Find $g(0)$ and solve $g(x) = 0$.

\item  Graph $g$.


\end{enumerate}

\item  Let $F = \{ (t, t^2) \, | \, \text{$t$ is a real number} \}$.  Find $F(4)$ and solve $F(x) = 4$.

\textbf{HINT:}  Elements of $F$ are of the form $(x, F(x))$.

\item  Let $G = \{ (2t, t+5) \, | \, \text{$t$ is a real number} \}$.  Find $G(4)$ and solve $G(x) = 4$.

\textbf{HINT:}  Elements of $G$ are of the form $(x, G(x))$.

%\setcounter{HW}{\value{enumi}}

\end{enumerate}



\begin{enumerate}
%\setcounter{enumi}{\value{HW}}

\item  The area enclosed by a square, in square inches,  is a function of the length of one of its sides $\ell$, when measured in inches.  This function is represented by the formula $A(\ell) = \ell^2$ for $\ell > 0$.  Find $A(3)$ and solve $A(\ell) = 36$.  Interpret your answers to each.  Why is $\ell$ restricted to $\ell > 0$?

\item  The area enclosed by a circle, in square meters, is a function of its radius $r$, when measured in meters.  This function is represented by the formula $A(r) = \pi r^2$ for $r > 0$.  Find $A(2)$ and solve $A(r) = 16\pi$.  Interpret your answers to each.  Why is $r$ restricted to $r > 0$?

\item  The volume  enclosed by a cube, in cubic centimeters, is a function of the length of one of its sides $s$, when measured in centimeters.   This function is represented by the formula $V(s) = s^3$ for $s > 0$.  Find $V(5)$ and solve $V(s) = 27$.  Interpret your answers to each.  Why is $s$ restricted to $s > 0$?

\item  The volume enclosed by a sphere, in cubic feet, is a function of the radius of the sphere $r$, when measured in feet. This function is represented by the formula  $V(r) =\frac{4\pi}{3} r^{3}$ for $r > 0$.  Find $V(3)$ and solve $V(r) = \frac{32\pi}{3}$.  Interpret your answers to each.  Why is $r$ restricted to $r > 0$?

\item  The height of an object dropped from the roof of an eight story building is modeled by the function:  $h(t) = -16t^2 + 64$, $0 \leq t \leq 2$. Here,  $h(t)$ is the height of the object off the ground, in feet, $t$ seconds after the object is dropped.  Find $h(0)$ and solve $h(t) = 0$.  Interpret your answers to each.  Why is $t$ restricted to $0 \leq t \leq 2$?

\item  The temperature in degrees Fahrenheit $t$ hours after 6 AM is given by $T(t) = -\frac{1}{2} t^2 + 8t+3$ for $0 \leq t \leq 12$. Find and interpret $T(0)$, $T(6)$ and $T(12)$.  

\item The function $C(x) = x^2-10x+27$  models the cost, in \textit{hundreds} of dollars, to produce $x$ \textit{thousand} pens.  Find and interpret $C(0)$, $C(2)$ and $C(5)$.

\item Using data from the  \href{http://www.bts.gov/publications/national_transportation_statistics/html/table_04_23.html}{\underline{Bureau of Transportation Statistics}}, the average fuel economy in miles per gallon for passenger cars in the US can be modeled by  $E(t) = -0.0076t^2+0.45t + 16$, $0 \leq t \leq 28$, where $t$ is the number of years since $1980$. Use a calculator to find $E(0)$, $E(14)$ and $E(28)$.  Round your answers to two decimal places and interpret your answers to each.

\item  The perimeter of a square, in centimeters,  is four times the length of one if its sides, also measured in centimeters.  Represent the function $P$ which takes as its input the length of the side of a square in centimeters, $s$ and returns the perimeter of the square in inches, $P(s)$ using a formula.

\item  The circumference of a circle, in feet,  is $\pi$ times the diameter of the circle, also measured in feet.  Represent the function $C$ which takes as its input the length of the diameter of a circle in feet, $D$ and returns the circumference of a circle in inches, $C(D)$ using a formula.

 

\item  Suppose $A(P)$ gives the amount of money in a retirement account (in dollars) after 30 years as a function of the amount of the monthly payment (in dollars), $P$.

\begin{enumerate}

\item What does $A(50)$ mean?  

\item  What is the significance of the solution to the equation $A(P) = 250000$? .

\item  Explain what each of the following expressions mean:  $A(P+50)$, $A(P)+50$, and $A(P) + A(50)$.  

\end{enumerate}

\item  Suppose $P(t)$ gives the chance of precipitation (in percent)  $t$ hours after 8 AM.

\begin{enumerate}

\item  Write an expression which gives the chance of precipitation at noon.

\item  Write an inequality which determines when the chance of precipitation is more than $50 \%$.

\end{enumerate}


\item Explain why the graph in  Exercise \ref{graphfunctionfirstvw}  suggests that not only is $v$ as a function of $w$ but also $w$ is a function of $v$.  Suppose $w = f(v)$ and $v = g(w)$.  That is, $f$ is the name of the function which takes $v$ values as inputs and returns $w$ values as outputs and $g$ is the name of the function which does vice-versa.   Find the domain and range of  $g$ and compare these to the domain and range of $f$.  

\item  Sketch the graph of a function with domain $(-\infty, 3) \cup [4,5)$ with range $\{ 2 \} \cup (5, \infty)$.

\end{enumerate}



 



\end{document}
